% !TeX program = XeLaTeX
% !TeX root = ../../shloka.tex

\sect{रामाष्टोत्तरशतनामस्तोत्रम्}

\dnsub{ध्यानम्}
\fourlineindentedshloka*
{श्रीराघवं दशरथात्मजमप्रमेयं}
{सीतापतिं रघुकुलान्वयरत्नदीपम्}
{आजानुबाहुमरविन्ददलायताक्षं}
{रामं निशाचरविनाशकरं नमामि}

\dnsub{स्तोत्रम्}
\twolineshloka
{श्रीरामो रामचन्द्रश्च रामभद्रश्च शाश्वतः}
{राजीवलोचनः श्रीमान् राजेन्द्रो रघुपुङ्गवः}

\twolineshloka
{जानकीवल्लभो जैत्रो जितामित्रो जनार्दनः}
{विश्वामित्रप्रियो दान्तः शरण्यत्राणतत्परः}

\twolineshloka
{वालिप्रमथनो वाग्मी सत्यवाक् सत्यविक्रमः}
{सत्यव्रतो व्रतधरः सदा हनुमदाश्रयः}

\twolineshloka
{कौसलेयः खरध्वंसी विराधवधपण्डितः}
{विभीषणपरित्राता हरकोदण्डखण्डनः}

\twolineshloka
{सप्ततालप्रभेत्ता च दशग्रीवशिरोहरः}
{जामदग्न्यमहादर्पदलनस्ताडकान्तकृत्}

\twolineshloka
{वेदान्तपारो वेदात्मा भवबन्धैकभेषजः}
{दूषणत्रिशिरोरिश्च त्रिमूर्तिस्त्रिगुणस्त्रयी}

\twolineshloka
{त्रिविक्रमस्त्रिलोकात्मा पुण्यचारित्रकीर्तनः}
{त्रिलोकरक्षको धन्वी दण्डकारण्यवासकृत्}

\twolineshloka
{अहल्यापावनश्चैव पितृभक्तो वरप्रदः}
{जितेन्द्रियो जितक्रोधो जितलोभो जगद्गुरुः}

\twolineshloka
{ऋक्षवानरसङ्घाती चित्रकूटसमाश्रयः}
{जयन्तत्राणवरदः सुमित्रापुत्रसेवितः}

\twolineshloka
{सर्वदेवाधिदेवश्च मृतवानरजीवनः}
{मायामारीचहन्ता च महाभोगो महाभुजः}

\twolineshloka
{सर्वदेवस्तुतः सौम्यो ब्रह्मण्यो मुनिसत्तमः}
{महायोगो महोदारः सुग्रीवस्थिरराज्यदः}

\twolineshloka
{सर्वपुण्याधिकफलः स्मृतः सर्वाघनाशनः}
{आदिपुरुषो महापुरुषः परमः पुरुषस्तथा}

\twolineshloka
{पुण्योदयो महासारः पुराणः पुरुषोत्तमः}
{स्मितवक्त्रो मितभाषी पूर्वभाषी च राघवः}

\twolineshloka
{अनन्तगुणगम्भीरो धीरो दान्तगुणोत्तमः}
{मायामानुषचारित्रो महादेवाभिपूजितः}

\twolineshloka
{सेतुकृज्जितवारीशः सर्वतीर्थमयो हरिः}
{श्यामाङ्गः सुन्दरः शूरः पीतवासा धनुर्धरः}

\twolineshloka
{सर्वयज्ञाधिपो यज्ञो जरामरणवर्जितः}
{शिवलिङ्गप्रतिष्ठाता सर्वाघगणवर्जितः}

\twolineshloka
{परमात्मा परं ब्रह्म सच्चिदानन्दविग्रहः}
{परञ्ज्योतिः परन्धाम पराकाशः परात्परः}

\twolineshloka
{परेशः पारगः पारः सर्वभूतात्मकः शिवः}
{इति श्रीरामचन्द्रस्य नाम्नामष्टोत्तरं शतम्}

\dnsub{फलश्रुतिः}
\resetShloka

\twolineshloka
{गुह्याद् गुह्यतरं देवि तव स्नेहात् प्रकीर्तितम्}
{यः पठेच्छृणुयाद् वाऽपि भक्तियुक्तेन चेतसा}

\twolineshloka
{स सर्वैर्मुच्यते पापैः कल्पकोटिशतोद्भवैः}
{जलानि स्थलतां यान्ति शत्रवो यान्ति मित्रताम्}

\twolineshloka
{राजानो दासतां यान्ति वह्नयो यान्ति सौम्यताम्}
{आनुकूल्यं च भूतानि स्थैर्यं यान्ति चलाः श्रियः}

\twolineshloka
{अनुग्रहं ग्रहा यान्ति शान्तिमायान्त्युपद्रवाः}
{पठतो भक्तिभावेन नरस्य गिरिसम्भवे}

\twolineshloka
{यः पठेत् परया भक्त्या तस्य वश्यं जगत्त्रयम्}
{यं यं कामं प्रकुरुते तं तमाप्नोति कीर्तनात्}

\twolineshloka
{कल्पकोटिसहस्राणि कल्पकोटिशतानि च}
{वैकुण्ठे मोदते नित्यं दशपूर्वैर्देशापरैः}

\twolineshloka
{रामं दूर्वादलश्यामं पद्माक्षं पीतवाससम्}
{स्तुवन्ति नामभिर्दिव्यैर्न ते संसारिणो जनाः}

\twolineshloka
{रामाय रामभद्राय रामचन्द्राय वेधसे}
{रघुनाथाय नाथाय सीतायाः पतये नमः}

\twolineshloka
{इमं मन्त्रं महादेवि जपन्नेव दिवानिशम्}
{सर्वपापविनिर्मुक्तो विष्णुसायुज्यमाप्नुयात्}


{॥इति~श्रीपद्मपुराणे~उत्तरखण्डे\\ श्रीरामाष्टोत्तरशतनामस्तोत्रं सम्पूर्णम्॥}
